\section{Objetivos}


\subsection{Objetivo General}

Desarrollar e implementar un sistema de navegación asistida en tiempo real para personas con discapacidad visual, basado en el uso de lentes adaptados con sensores de proximidad, visión artificial, trazado de rayos y tecnologías de inteligencia artificial, que permita identificar obstáculos, interpretar el entorno urbano y generar instrucciones auditivas contextuales. 

\subsection{Objetivos Específicos}

\begin{enumerate}
    
    \item Diseñar un modelo de visión por computadora que sea capaz de detectar y segmentar los objetos, identifique y delimite con precisión los objetos relevantes en el entorno de manera óptima y cuya precisión sea mayor 
al $80\%$.
    
    \item Reconstruir el escenario utilizando tecnologías relacionadas al raytracing para obtener una mayor fidelidad.
    
    \item Diseñar y desarrollar un modelo de navegación adecuado que permita al usuario desplazarse con autonomía y cuya medición del terreno sea en un $90\%$ segura, minimizando el riesgo de colisiones y caídas.  

    \item Modelar un sistema de interacción por voz que sea capaz de reconocer ciertas indicaciones del usuario y leer en voz alta las indicaciones del sistema.
    
    \item Desarrollar un modulo de lenguaje natural que traduzca y permita el manejo adecuado de indicaciones de manera cotidiana y durante las situaciones de emergencia.
\end{enumerate}

\section{Hipótesis}

Si se desarrolla un sistema de navegación asistida que integre visión por computadora, inteligencia artificial y guía auditiva en tiempo real, entonces será posible mejorar la seguridad y autonomía en el desplazamiento de personas con visión reducida al identificar obstáculos y proporcionar indicaciones de navegación durante su movilidad.

\subsection{Hipótesis específicas}
\begin{itemize}
    \item La utilización de modelos de detección de objetos basados en aprendizaje profundo permitirá identificar obstáculos relevantes en el entorno con suficiente precisión para aplicaciones de asistencia a la movilidad.

    \item La construcción de un mapa del entorno con localización simultánea permitirá generar rutas accesibles y proporcionar indicaciones auditivas al usuario durante su desplazamiento.

    \item La interacción mediante reconocimiento de voz facilitará el uso del sistema por parte de personas con discapacidad visual al eliminar la dependencia de interfaces visuales.
\end{itemize}